\documentclass[12pt,a4paper,openany]{book}

\usepackage[utf8]{inputenc}
\usepackage[T1]{fontenc}
\usepackage{lmodern}
\usepackage{microtype}
\usepackage{setspace}
\onehalfspacing
\emergencystretch=6em
\usepackage[margin=1in]{geometry}
\usepackage{amsmath,amssymb,amsthm,mathtools}
\usepackage{bm}
\usepackage{booktabs}
\usepackage{longtable}
\usepackage{multirow}
\usepackage{array}
\usepackage{graphicx}
\usepackage{float}
\usepackage{algorithm}
\usepackage{algpseudocode}
\usepackage{listings}
\usepackage{xcolor}
\usepackage{url}
\usepackage{hyperref}
\usepackage{cleveref}
\usepackage{natbib}
\usepackage{enumitem}
\usepackage{tikz}
\usetikzlibrary{positioning,arrows.meta,shapes.geometric}
\usepackage[most]{tcolorbox}

\hypersetup{
  colorlinks=true,
  linkcolor=blue!70!black,
  citecolor=green!40!black,
  urlcolor=blue!70!black
}

\definecolor{codebg}{gray}{0.96}
\definecolor{codegreen}{rgb}{0,0.45,0}
\definecolor{codegray}{rgb}{0.5,0.5,0.5}
\definecolor{codepurple}{rgb}{0.58,0,0.82}
\definecolor{passgreen}{rgb}{0.0,0.45,0.0}
\definecolor{failred}{rgb}{0.7,0.0,0.0}
\definecolor{warncolor}{rgb}{0.7,0.4,0.0}

\lstset{
  basicstyle=\ttfamily\small,
  keywordstyle=\color{blue!70!black}\bfseries,
  commentstyle=\color{codegreen},
  stringstyle=\color{codepurple},
  backgroundcolor=\color{codebg},
  showstringspaces=false,
  breaklines=true,
  frame=single,
  numbers=left,
  numberstyle=\tiny\color{codegray},
  numbersep=5pt
}

\theoremstyle{plain}
\newtheorem{theorem}{Theorem}[chapter]
\newtheorem{proposition}[theorem]{Proposition}
\newtheorem{lemma}[theorem]{Lemma}
\newtheorem{corollary}[theorem]{Corollary}

\theoremstyle{definition}
\newtheorem{definition}[theorem]{Definition}
\newtheorem{assumption}[theorem]{Assumption}
\newtheorem{example}[theorem]{Example}

\theoremstyle{remark}
\newtheorem{remark}[theorem]{Remark}

\newcommand{\E}{\mathbb{E}}
\newcommand{\R}{\mathbb{R}}
\newcommand{\N}{\mathbb{N}}
\newcommand{\calN}{\mathcal{N}}
\DeclareMathOperator{\Var}{Var}
\DeclareMathOperator{\Cov}{Cov}
\DeclareMathOperator{\diag}{diag}
\DeclareMathOperator*{\argmin}{arg\,min}
\DeclareMathOperator*{\argmax}{arg\,max}
\newcommand{\norm}[1]{\lVert #1 \rVert}
\DeclarePairedDelimiter{\abs}{\lvert}{\rvert}
\newcommand{\inner}[2]{\langle #1,\,#2 \rangle}
\newcommand{\kron}{\otimes}

\newtcolorbox{proposalbox}[1][]{
  colback=blue!3,
  colframe=blue!40!black,
  boxrule=0.6pt,
  arc=2pt,
  left=6pt,
  right=6pt,
  top=6pt,
  bottom=6pt,
  #1
}

\newtcolorbox{riskbox}[1][]{
  colback=orange!4,
  colframe=warncolor,
  boxrule=0.6pt,
  arc=2pt,
  left=6pt,
  right=6pt,
  top=6pt,
  bottom=6pt,
  #1
}


\title{\Huge\textbf{MathDevMCP}\\[10pt]
  \Large Final Release Report for an Industrial Mathematical Development Agent}
\author{Chak Wong and collaborators}
\date{April 29, 2026}

\begin{document}

\frontmatter
\maketitle
\tableofcontents

\mainmatter

\part{Release Claim}

\chapter{Executive Summary}

MathDevMCP is now a working internal release candidate for document-grounded
mathematical development. It provides a Python library, command-line tools, MCP
tool wrappers, benchmark gates, release profiles, backend diagnostics, and
privacy-preserving private-corpus validation. The product is designed for
colleagues who work across long LaTeX documents, mathematical code, proofs,
derivations, algorithms, and review workflows where unsupported claims are
expensive.

The release claim is deliberately narrow and evidence based. MathDevMCP does
not claim to prove arbitrary mathematical papers correct. It gives agents and
humans structured access to local documents, labels, code, parser evidence,
typed obligations, AST operation evidence, backend diagnostics, and release
gates. It helps colleagues find the right mathematical context, compare
documented claims to implementations, detect missing implementation terms,
route proof obligations, and abstain when evidence is not certifying.

The current full release profile is configured to require benchmark evidence,
parser policy, governance validation, release-corpus metadata, LeanDojo backend
environment evidence, LaTeXML validation, and a private-corpus manifest. The
private-corpus evidence used for this report is an external sanitized corpus
generated outside the repository. Its purpose is to validate the privacy and
release-gate machinery without committing private department documents. A real
department manifest can be substituted through the same environment variable.

\begin{proposalbox}
\textbf{Release status.} With an external sanitized private manifest supplied,
the full profile has no blocking findings. During report generation the only
remaining caveat is the dirty working tree caused by ongoing release edits.
\end{proposalbox}

\chapter{Product Scope}

MathDevMCP supports five colleague-facing activities:

\begin{enumerate}[label=\arabic*.]
  \item locating mathematical context in LaTeX source trees;
  \item extracting labeled equation, theorem, assumption, and paragraph
        neighborhoods with provenance;
  \item comparing document claims with Python implementations;
  \item auditing derivation obligations with parser, typed, shape, numeric, and
        backend evidence;
  \item producing repeatable release evidence across base, backend, LaTeXML,
        private-corpus, and full profiles.
\end{enumerate}

The product is intentionally local-first. Source documents and code remain on
the user's machine. Optional tools such as LaTeXML, Lean, LeanDojo, Sage, and
SymPy are detected at runtime and reported through structured diagnostics.
LeanDojo remains in the isolated \texttt{mathdevmcp-backends} environment so its
dependency stack does not destabilize the main working environment.

\section{What is included}

The release includes:

\begin{itemize}
  \item a LaTeX indexer with label, environment, and context extraction;
  \item code and document search utilities;
  \item code--document consistency checks;
  \item derivation-step checks with explicit \emph{unverified} and
        \emph{mismatch} outcomes;
  \item proof obligation and proof-audit workflows;
  \item typed and dimensional diagnostic workflows;
  \item AST operation extraction for mathematical implementation patterns;
  \item parser benchmark and parser policy selection;
  \item release corpus manifest validation;
  \item private-corpus validation with redacted output;
  \item release-readiness profiles;
  \item MCP facade and FastMCP server wrappers;
  \item operator, deployment, release, security, and maintainer documentation.
\end{itemize}

\section{What is not claimed}

MathDevMCP does not replace expert mathematical review. A parser hit is not a
proof. A shape diagnostic is not a theorem. A Lean skeleton is not a certificate
until direct Lean checking accepts it without placeholders. A numeric diagnostic
can refute a false claim or suggest a failure mode, but it does not establish a
general theorem. The product's value is that it preserves these distinctions in
machine-readable reports.

\chapter{Release Profiles}

The release profiles are:

\begin{description}
  \item[base] Public benchmarks, parser policy, governance, release corpus, and
    normal doctor evidence.
  \item[backend] Base evidence plus isolated LeanDojo backend environment
    validation.
  \item[latexml] Base evidence plus strict LaTeXML parser validation.
  \item[private-corpus] Base evidence plus an external private or sanitized
    private manifest.
  \item[full] All of the above.
\end{description}

The profile split is important because colleagues may use the base product
without every optional backend, while release managers can demand strict
evidence before a broader deployment. Optional evidence is not silently ignored:
strict profiles turn missing evidence into blocking findings.

\section{Full profile evidence}

\lstinputlisting[caption={Full profile release-readiness summary},basicstyle=\ttfamily\footnotesize]{generated/release_report/release-readiness-full-summary.txt}

\section{Doctor evidence}

\lstinputlisting[caption={Doctor capability summary},basicstyle=\ttfamily\footnotesize]{generated/release_report/doctor-summary.txt}

\part{Architecture}

\chapter{System Architecture}

MathDevMCP has three layers. The first layer is a Python implementation
substrate that owns parsing, indexing, consistency checks, proof-audit routing,
benchmarking, release policy, and diagnostics. The second layer is a CLI and MCP
facade that presents stable tools to users and agents. The third layer is a
workflow and documentation layer that explains how to combine tools for
colleague tasks.

\section{Implementation substrate}

The Python package is organized around small report-producing functions. Each
public report includes a metadata contract so that CLI, MCP, tests, and release
docs can depend on stable shapes. For example, release readiness returns a
\texttt{release\_readiness\_report}; parser benchmark returns
\texttt{parser\_backend\_result}; private corpus validation returns
\texttt{private\_corpus\_validation\_report}. This is the main compatibility
discipline for future maintainers.

\section{CLI and MCP facades}

The CLI is the reproducible interface for humans and CI. The MCP facade wraps
the same library functions so an agent can call them interactively. The facade
does not grant unrestricted shell access. It exposes bounded tools such as
\texttt{search\_latex}, \texttt{compare\_label\_code},
\texttt{audit\_derivation\_v2\_label}, \texttt{benchmark\_gate}, and
\texttt{release\_readiness}.

\section{Backend isolation}

Optional backends are treated as evidence engines. LaTeXML, Pandoc, Lean, Sage,
SymPy, and LeanDojo are detected by \texttt{doctor}. LeanDojo is checked through
the backend environment when configured. This avoids dependency conflicts in
the main Python environment while still allowing proof-search experiments and
Lean fixture validation.

\chapter{Core Data Contracts}

Every user-facing workflow should be understood as a contract. A contract has a
schema version, a name, a status or result payload, and where appropriate
provenance. This is why downstream documentation, agents, and tests can
interpret outputs without scraping prose.

\section{Status vocabulary}

The primary statuses are:

\begin{description}
  \item[consistent] The requested terms or structure were found in the checked
    context.
  \item[equivalent] A scoped derivation step matched exactly after normalization
    or backend checking.
  \item[unverified] Evidence exists, but it is not enough for certification.
  \item[mismatch] Required structure is missing or contradicted.
  \item[inconclusive] The tool lacks enough structure, parser support, or
    backend evidence to decide.
\end{description}

This vocabulary is central to the product. It prevents the agent from turning a
nearby equation or diagnostic hint into a categorical mathematical claim.

\chapter{Parser Policy}

MathDevMCP measures parser behavior before relying on parser output for proof
audit. The current lightweight parser is selected for proof audit when it
preserves expected labels and line provenance. LaTeXML is also validated in the
strict LaTeXML profile, but current release routing prefers the provenance and
environment structure available through the current parser.

\lstinputlisting[caption={Parser policy summary},basicstyle=\ttfamily\footnotesize]{generated/release_report/parser-policy-summary.txt}

\chapter{Benchmark Gate}

The benchmark gate is the main regression control. It covers consistency,
derivation, workflow, proof-audit, parser-corpus, AST-corpus, typed IR,
industrial review, and release-corpus cases. The gate currently requires every
benchmark case to pass; expected abstentions are counted as quality signals, not
as failures.

\lstinputlisting[caption={Benchmark gate summary},basicstyle=\ttfamily\footnotesize]{generated/release_report/benchmark-gate-summary.txt}

\part{Colleague Workflows}

\chapter{Workflow 1: Find Mathematical Context}

A colleague begins with a topic such as a likelihood, Hamiltonian flow, Euler
equation, or stability condition. The first tool is usually
\texttt{search-latex}. It returns document blocks with labels and provenance so
the agent can ground a response in the source document.

\lstinputlisting[caption={Search workflow output},basicstyle=\ttfamily\scriptsize]{generated/release_report/workflow-search.txt}

\chapter{Workflow 2: Read a Labeled Neighborhood}

Once a label is known, the user can extract a paragraph neighborhood. This is
more useful than a raw line excerpt when an agent needs nearby explanatory
prose, assumptions, or notation.

\lstinputlisting[caption={Labeled neighborhood output},basicstyle=\ttfamily\scriptsize]{generated/release_report/workflow-context.txt}

\chapter{Workflow 3: Compare Document and Code}

Code--document drift is one of the highest-value release use cases. The
\texttt{compare-label-code} workflow checks whether required terms from a
labeled document block are present in a code file. It is not a full proof of
equivalence. It is a targeted drift detector.

\lstinputlisting[caption={Code--document comparison output},basicstyle=\ttfamily\scriptsize]{generated/release_report/workflow-compare.txt}

\chapter{Workflow 4: Audit a Derivation}

The v2 audit workflow combines parser policy, typed obligations, route
decisions, shape diagnostics, numeric suggestions, backend evidence, and
verification-boundary text. It is the most appropriate colleague-facing surface
when a user asks whether a labeled mathematical statement is actually certified.

\lstinputlisting[caption={Derivation audit output},basicstyle=\ttfamily\scriptsize]{generated/release_report/workflow-audit.txt}

\chapter{Workflow 5: Build an Implementation Brief}

The implementation brief bundles search, label selection, context extraction,
consistency checking, and optional derivation checking into one object. This is
the recommended tool when another agent needs a document-grounded starting
point for code review or implementation.

\lstinputlisting[caption={Implementation brief output},basicstyle=\ttfamily\scriptsize]{generated/release_report/workflow-brief.txt}

\part{Case Studies}

\chapter{Kalman State-Space Likelihood}

The Kalman likelihood case checks whether implementation code carries the
linear algebra structure described by the document. Missing solves, inverses,
or log-determinant terms are important because they often indicate stale
documentation or a numerically fragile implementation. MathDevMCP reports such
cases as consistency findings rather than silently rationalizing them.

For a colleague, the useful behavior is fast localization. The tool identifies
the label, extracts the surrounding text, compares the implementation file, and
reports missing required terms. The next action is human review of the
likelihood implementation and associated tests.

\chapter{HMC Leapfrog and Hamiltonian Flow}

Hamiltonian Monte Carlo code is sensitive to sign conventions, gradients,
Metropolis corrections, and leapfrog ordering. MathDevMCP does not certify
convergence. It checks that documented Hamiltonian and leapfrog structures are
visible in the implementation and routes uncertified claims to abstention when
regularity or ergodicity assumptions are not present.

\chapter{Macro Filter Multi-File Corpus}

The macro filter fixture exercises multi-file LaTeX inclusion, macro
definitions, repeated notation, and a deliberately missing gain update. This is
close to real colleague work: a main document includes a model file, macros are
defined in one place, and code may omit a step that the exposition assumes.

\chapter{DSGE Euler Equation}

The DSGE case checks Euler residual and stochastic discount factor structure.
MathDevMCP helps a reviewer avoid claiming that a model is solved or calibrated
unless the assumptions and implementation evidence are available. The expected
abstention is part of the release quality signal.

\chapter{Stochastic Volatility Likelihood}

The stochastic volatility case checks transition and likelihood labels. The
important release behavior is conservative handling of latent-volatility prior
assumptions and Jacobian-like omissions. A diagnostic may expose a missing
term, but the tool should not overstate proof.

\chapter{SDE and PDE Numerics}

The SDE/PDE numerics case covers Euler--Maruyama and residual/stability
structure. A colleague can use the tool to locate the discretization formula,
check whether code carries time-step and stability terms, and record a
discretization-assumption review when the evidence is incomplete.

\chapter{ML and LLM Objective Functions}

The ML objective case catches wrong-sign gradients and missing objective terms.
This is useful for training and fine-tuning code where the prose can describe a
loss while the implementation evolves quickly. MathDevMCP provides a structured
way to ask whether the documented loss and gradient are visible in code.

\chapter{Bayesian ELBO and Variational Inference}

The ELBO case checks objective, expectation, entropy, and reparameterization
structure. The release behavior emphasizes abstention around variational-family
assumptions and support conditions. This keeps the agent from treating a
surface-level match as a full variational proof.

\chapter{Computational Physics MCMC}

The computational physics case checks acceptance ratios, Hamiltonian energy,
gradients, and leapfrog-like operations. The tool can surface a missing
Metropolis correction or energy term, but it deliberately leaves convergence
and ergodicity claims for expert review.

\chapter{Private Corpus Validation}

The private-corpus workflow validates external manifests without committing
private material. Normal output redacts private paths. The generated evidence
below uses a sanitized external corpus under the same release machinery that a
real private department manifest would use.

\lstinputlisting[caption={Private corpus validation summary},basicstyle=\ttfamily\footnotesize]{generated/release_report/private-corpus-summary.txt}

\part{Security, Operations, and Maintenance}

\chapter{Security and Privacy}

Private documents must stay outside the repository. A populated private
manifest may contain sensitive paths and should remain external. Normal
MathDevMCP reports redact private document roots and code roots. Governance
validation rejects private paths inside the checkout. Evidence bundles for
review should be stored outside git unless they are deliberately reduced to
safe report snippets.

\chapter{Operations}

The normal operator sequence is:

\begin{lstlisting}[language=bash]
PYTHONPATH=src python -m mathdevmcp.cli doctor
scripts/backend_env_doctor.sh "$PWD"
MATHDEVMCP_REQUIRE_LATEXML=1 scripts/validate_latexml_backend.sh "$PWD"
PYTHONPATH=src python -m mathdevmcp.cli benchmark-gate --root "$PWD"
PYTHONPATH=src python -m mathdevmcp.cli release-readiness --root "$PWD" --profile base
\end{lstlisting}

For full release validation, add:

\begin{lstlisting}[language=bash]
export MATHDEVMCP_PRIVATE_CORPUS_MANIFEST=/secure/local/path/corpus.json
scripts/validate_private_corpus.sh "$PWD"
PYTHONPATH=src python -m mathdevmcp.cli release-readiness --root "$PWD" --profile full
\end{lstlisting}

\chapter{Maintainer Guide}

Maintain release behavior by changing library code first, then CLI/MCP wrappers,
then docs and evidence. When adding a new benchmark category, update the case
builder, runner, summary tests, release report evidence, and release policy if
the category becomes mandatory. When adding a new private corpus entry, keep the
entry outside git and check redaction before sharing output.

The highest-risk modules are:

\begin{itemize}
  \item \texttt{benchmarks.py}, for benchmark totals and gate behavior;
  \item \texttt{release\_policy.py}, for profile blockers and caveats;
  \item \texttt{release\_corpus.py}, for privacy redaction and manifest shape;
  \item \texttt{parser\_policy.py}, for proof-audit parser routing;
  \item \texttt{proof\_audit\_v2.py}, for verification-boundary reporting;
  \item \texttt{mcp\_facade.py} and \texttt{cli.py}, for public tool surfaces.
\end{itemize}

\chapter{Limitations and Accepted Boundaries}

The product is conservative by design. It is useful because it knows when to
stop. It does not convert every mathematical statement to a theorem prover. It
does not guarantee semantic equivalence between arbitrary prose and code. It
does not make private data safe to commit. It does not make numeric diagnostics
into proofs. These boundaries should remain explicit in future releases.

\part{Detailed Release Manual}

\chapter{End-to-End Colleague Scenario}

This chapter walks through a realistic review day. A colleague has a document
claim about a state-space likelihood and a Python implementation that is being
edited quickly before a deadline. The colleague does not want a polished essay;
they want to know where the claim lives, whether the code visibly contains the
required numerical operations, whether a derivation can be certified, and what
should be reviewed by a human.

\section{Step 1: locate the candidate claim}

The agent begins with \texttt{search-latex}. This avoids relying on a language
model's memory of the repository. The search result gives labels and local
context. The colleague can inspect those labels directly or ask the agent to use
the top result. If no label is found, the workflow stops as \emph{inconclusive}
rather than inventing a source.

\section{Step 2: extract neighborhood context}

The label neighborhood gives surrounding prose and assumptions. This step
matters because many mathematical claims are not self-contained. A likelihood
equation may rely on notation introduced one paragraph above, while a stability
condition may depend on a boundary condition or discretization convention.

\section{Step 3: compare to code}

The comparison step asks for concrete required terms such as
\texttt{logdet}, \texttt{solve}, \texttt{gradient}, or \texttt{stability}. A
missing term is not automatically a proof of a bug, but it is a precise review
signal. The output includes the selected label, findings, and status, so the
agent can explain what was checked.

\section{Step 4: audit the derivation boundary}

The v2 audit report is where MathDevMCP becomes most valuable for disciplined
agent work. It records parser policy, typed diagnostics, shape diagnostics,
numeric diagnostic suggestions, backend attempts, and a verification boundary.
That boundary text is essential: it tells the colleague whether the output is a
certificate, a diagnostic, or an abstention.

\section{Step 5: produce an implementation brief}

For handoff, the implementation brief bundles the key evidence. A teammate can
read one object rather than rerun every command. This is especially useful when
the agent is asked to create a code-review note, a pull-request comment, or a
maintenance ticket.

\chapter{Release Gate Interpretation}

The release gate should be read as a contract, not a vibe check. A green base
profile means public fixtures, governance, release corpus metadata, parser
policy, and benchmark cases pass. It does not mean every optional backend or
private corpus is available. A green full profile means the strict optional
evidence has also been supplied. This separation lets the team use MathDevMCP
productively on ordinary machines while still preserving a stricter bar for
industrial release.

\section{Blockers}

Blockers represent release conditions that must be fixed or explicitly scoped
out. Examples include benchmark failure, parser policy failure, governance
validation failure, missing strict LaTeXML evidence in the LaTeXML profile,
missing backend evidence in the backend profile, and missing private manifest
evidence in the private-corpus or full profiles.

\section{Caveats}

Caveats are not ignored. They are tracked because a release may be useful while
still carrying environmental or evidence limitations. A dirty worktree is a
caveat because generated reports should name the actual commit. Missing private
corpus evidence is a low-severity caveat in the base profile, but a blocker in
the private-corpus and full profiles.

\section{Expected abstentions}

Expected abstentions protect the product from false confidence. A benchmark can
pass because the tool correctly abstained in a case where certification would
be inappropriate. This is different from a tolerated failure budget. The
current benchmark policy requires all cases to pass.

\chapter{Private Corpus Operating Model}

Private corpus validation is designed for teams that cannot commit research
documents, code paths, or real local directory structures. The populated
manifest lives outside git. The release validator reads it locally, uses
unredacted paths internally to parse documents, and emits redacted reports.

\section{External sanitized corpus}

For autonomous release validation in this workspace, the execution created an
external sanitized corpus under a temporary directory. It includes six release
domains: DSGE macro-finance, stochastic volatility, SDE/PDE numerics, ML/LLM
objectives, Bayesian ELBO/VI, and computational-physics MCMC. The generated
manifest is not committed. The report includes only redacted summaries.

\section{Real private corpus substitution}

A real deployment should replace the sanitized manifest with an approved
external manifest. The same validator and full profile commands apply. The
operator should verify that every release-gated entry has expected labels,
parser backends, expected abstentions or false-confidence seeds, and external
document roots. If a real private corpus covers fewer domains than the
sanitized release corpus, the release report should say exactly that.

\section{Failure modes}

Important private-corpus failure modes are:

\begin{itemize}
  \item missing manifest path;
  \item invalid JSON;
  \item malformed entry fields;
  \item private path inside the repository checkout;
  \item missing document root or code root;
  \item release-gated entry with no expected labels;
  \item release-gated entry with no parser backend;
  \item parser policy not selected for proof audit.
\end{itemize}

\chapter{Backend Environment Operating Model}

Backend tools are intentionally not all installed into the main Python
environment. This is a release-quality design choice. LeanDojo can bring a
dependency stack that conflicts with unrelated document or ML packages. The
backend env checks therefore prove that the optional backend can be invoked
without requiring the main working environment to import every optional module.

\section{Lean}

Lean is detected as an executable, usually through the user-local
\texttt{elan} proxy. The release pins the intended toolchain through
\texttt{MATHDEVMCP\_LEAN\_TOOLCHAIN} for MathDevMCP subprocesses. This avoids
changing the user's global Lean default.

\section{LeanDojo}

LeanDojo is checked as a Python module inside \texttt{mathdevmcp-backends}. A
plain \texttt{doctor} command in the active environment may report LeanDojo as
unavailable. The backend env doctor is the authoritative check for the backend
profile.

\section{LaTeXML}

LaTeXML is a system executable, not a Python dependency. The strict LaTeXML
profile requires it to preserve expected labels on the benchmark corpus. The
current evidence shows label preservation and source provenance. Environment
recognition differs from the lightweight parser, so parser policy still chooses
the current parser for proof-audit routing.

\chapter{Code Architecture for Maintainers}

The implementation should be maintained by preserving the boundary between
library logic, CLI rendering, MCP wrappers, scripts, and documentation.

\section{Library logic}

Library functions build structured dictionaries with contract metadata. They
should be deterministic for the same input files and environment. New behavior
should be added in library code first and exercised by direct tests.

\section{CLI}

The CLI should remain a thin argument parser. It should call library functions,
print JSON, and return meaningful exit codes. It should not contain hidden
business logic that bypasses the tested library surface.

\section{MCP facade}

The MCP facade should wrap the same library functions. Tool names are part of
the colleague-facing API. If a tool is renamed, keep a compatibility alias or
document a breaking change.

\section{Scripts}

Scripts should orchestrate repeatable workflows. They should set
\texttt{PYTHONPATH} explicitly, avoid shell string evaluation for backend
commands, and preserve redaction. Scripts that create generated evidence should
write to ignored or explicitly requested output directories.

\chapter{Adding a New Domain}

To add a new mathematical domain, follow this release path:

\begin{enumerate}[label=\arabic*.]
  \item Add a small public or sanitized fixture with at least one labeled
        mathematical block.
  \item Add paired code showing the relevant operations or a seeded mismatch.
  \item Extend AST operation extraction only if the new operation is genuinely
        needed.
  \item Add benchmark cases for parser preservation, AST operation coverage,
        and expected abstention behavior.
  \item Add release corpus metadata.
  \item Add a report case study and generated evidence snippet if the domain is
        release-facing.
  \item Update private manifest templates only with placeholder paths.
\end{enumerate}

This path keeps new domains measurable. It also prevents the product from
accumulating impressive-sounding workflows that do not have regression tests.

\chapter{Adding a New Backend}

A new backend should enter the system as optional evidence. The correct path is
to add a doctor capability, an isolated invocation strategy if needed, timeout
protection, a clear result contract, and tests for unavailable, timeout,
inconclusive, mismatch, and certifying behavior. Only after that should the
backend be used in proof-audit routing or release profiles.

\section{Backend result discipline}

Backend results should distinguish:

\begin{itemize}
  \item backend unavailable;
  \item input not encodable;
  \item timeout;
  \item diagnostic-only evidence;
  \item counterexample or mismatch;
  \item deterministic certificate.
\end{itemize}

The last item is intentionally rare. Most mathematical development tools are
useful because they reduce search space and expose gaps, not because they
automatically certify every claim.

\chapter{Release Evidence Maintenance}

The evidence snippets in this report are generated by
\texttt{scripts/generate\_release\_report\_evidence.sh}. Regenerate them after
changing release policy, parser policy, benchmark cases, private corpus
validation, or user-facing workflow examples. Do not hand-edit generated
evidence unless the file is explicitly converted into narrative documentation.

\section{Evidence hygiene}

After regeneration, run a path leak scan. The scan should check for real private
roots, local manifest filenames, and repository-local absolute paths. The
committed snippets may contain \texttt{<repo>},
\texttt{<redacted-private-path>}, and
\texttt{<redacted-private-manifest>}.

\section{Evidence reproducibility}

Generated snippets should name the command that produced them. A colleague
reading the report should be able to rerun the command in the repository and
understand why the local result may differ if the worktree is dirty, optional
backends are missing, or a private manifest is not configured.

\chapter{Risk Register}

\section{Risk: overclaiming certification}

The most important product risk is false confidence. Mitigation: keep
verification-boundary text in proof-audit reports, preserve diagnostic-only
flags, and require benchmark cases where abstention is the correct outcome.

\section{Risk: private data leakage}

Private roots and manifests are sensitive. Mitigation: keep populated manifests
outside git, redact normal reports, reject private paths inside the checkout,
and scan generated docs before commit.

\section{Risk: backend dependency conflict}

Optional backends can conflict with the main environment. Mitigation: use
\texttt{mathdevmcp-backends}, backend-specific doctor checks, and explicit
environment variables rather than broad installation into the base env.

\section{Risk: benchmark drift}

Benchmark totals and category meanings can drift as the project grows.
Mitigation: update tests, reset memo, and release report together whenever a
benchmark case is added or removed.

\section{Risk: large-module entropy}

The codebase has several large modules. Mitigation: keep public facades stable,
extract pure helper modules when it improves locality, and run focused tests
after each extraction.

\appendix
\part*{Appendices}
\addcontentsline{toc}{part}{Appendices}

\chapter{Release Corpus Summary}

\lstinputlisting[caption={Release corpus validation summary},basicstyle=\ttfamily\footnotesize]{generated/release_report/release-corpus-summary.txt}

\chapter{Benchmark Gate JSON Excerpt}

\lstinputlisting[caption={Benchmark gate JSON excerpt},basicstyle=\ttfamily\scriptsize]{generated/release_report/benchmark-gate-json-excerpt.txt}

\chapter{Parser Policy JSON Excerpt}

\lstinputlisting[caption={Parser policy JSON excerpt},basicstyle=\ttfamily\scriptsize]{generated/release_report/parser-policy-json-excerpt.txt}

\chapter{Release Corpus JSON Excerpt}

\lstinputlisting[caption={Release corpus JSON excerpt},basicstyle=\ttfamily\scriptsize]{generated/release_report/release-corpus-json-excerpt.txt}

\chapter{Full Release Readiness JSON Excerpt}

\lstinputlisting[caption={Full release readiness JSON excerpt},basicstyle=\ttfamily\scriptsize]{generated/release_report/release-readiness-full-json-excerpt.txt}

\chapter{Private Corpus JSON Excerpt}

\lstinputlisting[caption={Private corpus JSON excerpt},basicstyle=\ttfamily\scriptsize]{generated/release_report/private-corpus-json-excerpt.txt}

\chapter{Evidence Index}

\lstinputlisting[caption={Generated evidence index},basicstyle=\ttfamily\footnotesize]{generated/release_report/evidence-index.txt}

\chapter{Command Reference}

\begin{longtable}{p{0.32\textwidth}p{0.58\textwidth}}
\toprule
Command & Purpose \\
\midrule
\texttt{doctor} & Report runtime capabilities and optional backend status. \\
\texttt{search-latex} & Search a LaTeX source tree for relevant context. \\
\texttt{extract-latex-neighborhood} & Extract paragraph context around a label. \\
\texttt{compare-label-code} & Compare labeled document terms to code. \\
\texttt{audit-derivation-v2-label} & Build typed proof-audit evidence for a label. \\
\texttt{implementation-brief} & Bundle search, context, and code checks. \\
\texttt{benchmark-gate} & Run the release benchmark gate. \\
\texttt{validate-release-corpus} & Validate public and configured private corpus metadata. \\
\texttt{release-readiness} & Produce a profile-specific release report. \\
\bottomrule
\end{longtable}

\chapter{Private Manifest Fields}

\begin{longtable}{p{0.32\textwidth}p{0.58\textwidth}}
\toprule
Field & Meaning \\
\midrule
\texttt{id} & Stable entry identifier. \\
\texttt{domain} & Mathematical or colleague workflow domain. \\
\texttt{privacy\_class} & \texttt{private\_external} or \texttt{private\_sanitized\_external}. \\
\texttt{document\_root} & External LaTeX/document root. \\
\texttt{code\_roots} & External code roots associated with the documents. \\
\texttt{expected\_labels} & Labels the parser must preserve. \\
\texttt{expected\_operations} & Mathematical operations expected in code or AST evidence. \\
\texttt{expected\_abstentions} & Known cases where conservative abstention is expected. \\
\texttt{seeded\_false\_confidence\_cases} & False-confidence seeds for review. \\
\texttt{required\_parser\_backends} & Parser backends required for the entry. \\
\texttt{release\_gate\_enabled} & Whether the entry is mandatory for private profile evidence. \\
\texttt{notes} & Human-readable release note. \\
\bottomrule
\end{longtable}

\backmatter
\bibliographystyle{plainnat}
\bibliography{references}

\end{document}
